% Pass: Opus 5 (claude-opus-5[1m]), 2026-09-07 - first pass, unchecked against
% the sheets by a human.
%
% Dealers/Etchings, batch 4: twelve sheets, leaves 88 to 98 and the back board.
% This is the last batch of the volume and the last of the six. Eleven of the
% twelve sheets carry writing; the twelfth is the back board, which carries
% none. No opening is photographed twice, and - as in batches 1, 2 and 3 -
% there is not one clipping in the batch.
%
% Four things about these twelve leaves.
%
% First, the printed numbers run 88, 89, 90 ... 98 without a gap. That has not
% happened before in this volume: batch 2 lost 64 and 65, batch 3 lost 79. The
% book ends where the paper ends.
%
% Second, leaf 51's index of dealers is spent by the foot of leaf 89. Leaves 88
% and 89 clear the last nine accounts - Mrs. Albert Sterner, E. & A. Milch, the
% American Society of Graphic Arts, Sidney Phillips, Babcock, the Print Rooms
% of San Francisco, Carl Smalley of MacPherson Kansas, Brown Robertson, and
% Robertson Deschamps - most of them one small consignment of etchings at
% eighteen dollars, most of them returned. From leaf 90 on there is one dealer
% and it is Frank K. M. Rehn, whose account runs across leaves 90, 91, 92, 93,
% 94 and 95 and is the whole back of the book.
%
% Third, leaves 92 and 93 are not an account at all. They are the list of what
% went to Rehn from Gloucester and Rockland in August and September 1926, and
% each picture carries a sentence: « black smoke stack forward, shore & grey
% sky », « derelict drawn ashore, gloomy looking », « old boarded up boarding
% house, Rockland, near shipyard », « wild & wooly ». Two of the thirty were
% bought by Rehn himself in October at 200 less a third. In the left margin she
% goes on judging the pictures after the fact - « Painted spring 1926, brought
% to R. later in fall - not so good ».
%
% Fourth, the volume closes on her own work. « J. H. in rear studio, 3 Wash.
% Sq. N. » heads leaf 96, and leaves 96, 97 and 98 inventory Josephine Nivison
% Hopper's canvases, watercolours, pastels and little oils, then where they
% were shown and what they were priced at: the Whitney Club in 1927, the Morton
% Gallery in 1928, the Corcoran Biennial in 1951, and « Mar. 18" '46. Turned
% down by jury of [torn] Club (thy lowest depths [torn] Pen & Brush! » She
% wrote every account in the six volumes; this is the only part of them about
% her. The last line of the book is a watercolour shown in a parish house in
% Lent 1952, and the fore-edge of that leaf is torn away.
\documentclass[11pt,a4paper]{article}
% The preamble shared by every transcription of the Hopper artist's ledgers.
%
% It rests on one idea, taken whole from the Grothendieck workbench this
% method comes from: **uncertainty compounds, it does not resolve.** A
% transcription that smooths over a doubtful figure has destroyed the only
% thing it was made for — a reader must be able to tell, at every point, what
% was on the leaf and what was supplied.
%
% A ledger sharpens that. In a manuscript of mathematics an invented word is a
% wrong word. Here an invented figure is a *sale that did not happen*, at a
% price nobody paid, to a buyer who never bought — and it will be cited,
% because a table looks like data in a way that prose does not.
%
% The same commands are recognised by `scripts/render.mjs`, which produces the
% site's reading view. A macro added here must be added there too, or rendering
% fails loudly, which is the wanted behaviour: a silently wrong rendering is
% worse than none.

% \ProvidesFile, not \ProvidesPackage: the transcriptions pull this in with
% % The preamble shared by every transcription of the Hopper artist's ledgers.
%
% It rests on one idea, taken whole from the Grothendieck workbench this
% method comes from: **uncertainty compounds, it does not resolve.** A
% transcription that smooths over a doubtful figure has destroyed the only
% thing it was made for — a reader must be able to tell, at every point, what
% was on the leaf and what was supplied.
%
% A ledger sharpens that. In a manuscript of mathematics an invented word is a
% wrong word. Here an invented figure is a *sale that did not happen*, at a
% price nobody paid, to a buyer who never bought — and it will be cited,
% because a table looks like data in a way that prose does not.
%
% The same commands are recognised by `scripts/render.mjs`, which produces the
% site's reading view. A macro added here must be added there too, or rendering
% fails loudly, which is the wanted behaviour: a silently wrong rendering is
% worse than none.

% \ProvidesFile, not \ProvidesPackage: the transcriptions pull this in with
% % The preamble shared by every transcription of the Hopper artist's ledgers.
%
% It rests on one idea, taken whole from the Grothendieck workbench this
% method comes from: **uncertainty compounds, it does not resolve.** A
% transcription that smooths over a doubtful figure has destroyed the only
% thing it was made for — a reader must be able to tell, at every point, what
% was on the leaf and what was supplied.
%
% A ledger sharpens that. In a manuscript of mathematics an invented word is a
% wrong word. Here an invented figure is a *sale that did not happen*, at a
% price nobody paid, to a buyer who never bought — and it will be cited,
% because a table looks like data in a way that prose does not.
%
% The same commands are recognised by `scripts/render.mjs`, which produces the
% site's reading view. A macro added here must be added there too, or rendering
% fails loudly, which is the wanted behaviour: a silently wrong rendering is
% worse than none.

% \ProvidesFile, not \ProvidesPackage: the transcriptions pull this in with
% \input{../preamble/hopper} rather than \usepackage, and \ProvidesPackage in
% an \input-ed file makes LaTeX warn on every compile that it "requested
% package `'" and got `hopper'. A warning nobody can act on is a warning
% everybody learns to scroll past, which is how the real ones get missed.
\ProvidesFile{hopper.sty}[2026/09/05 Transcriptions of the Hopper artist's ledgers]

\usepackage{array}
% longtable, not tabularx: a leaf of thirty-one exhibition lines must break
% across pages, and tabularx cannot be wrapped in an environment of our own —
% it scans the source for a literal \end{tabularx} and dies on \end{ledgertable}
% with « File ended while scanning use of \TX@get@body ». longtable has no such
% scan, and the wrapping column type below does the job tabularx's X would.
\usepackage{longtable}

% A wrapping column, `Y{0.3}` being three tenths of the measure.
%
% Prose columns must be declared this way and not as `l`. A table of `l`
% columns is set to its natural width, which can be wider than the page, and
% TeX issues **no warning at all** because nothing ever asked the row to fit:
% the first compile of Book I produced a page whose right-hand column ran off
% the paper with a clean log. A fixed-width column wraps, and if it still
% cannot fit, TeX says so — which is what `npm run pdf` then refuses to build
% over.
%
% Leave about 0.03 of the measure per column for the inter-column space: five
% columns can share 0.85, six can share 0.80. Getting it wrong is not a
% guessing game — `npm run pdf` reports the overfull box and how many points
% too wide it is.
\newcolumntype{Y}[1]{>{\raggedright\arraybackslash}p{#1\linewidth}}
\usepackage[english]{babel}
\usepackage{fontspec}
\usepackage{geometry}
\geometry{a4paper,margin=25mm,marginparwidth=28mm}
\usepackage{marginnote}
\usepackage{xcolor}
\usepackage{tikz}
% ulem, not soul. `soul`'s \st and \ul are built on a character-by-character
% scan that XeLaTeX's font machinery breaks: the document fails with an
% undefined control sequence rather than degrading. `normalem` stops it
% hijacking \emph, which the transcriptions use for underlining of Jo Hopper's
% own.
\usepackage[normalem]{ulem}
% ulem sets each underlined word in a box TeX can neither hyphenate nor break
% after, so a paragraph dense in \uncertain{} readings can reach a state where
% no legal breakpoint exists. \emergencystretch lets TeX loosen it on a third
% pass instead of overrunning: an overfull line in a transcription hides
% content.
\setlength{\emergencystretch}{3em}
\usepackage{hyperref}
\hypersetup{colorlinks=true,linkcolor=black,urlcolor=black,pdfborder={0 0 0}}

% --- Batch metadata ---------------------------------------------------------
% Taken as they stand by the rendering script, which shows them at the head of
% the reading view. They fix what the file claims to cover: without them a
% transcription floats unmoored in an archive of five hundred sheets.

\newcommand{\ledger}[1]{\gdef\hopLedger{#1}}
\newcommand{\ledgertitle}[1]{\gdef\hopLedgerTitle{#1}}
% A notebook of Josephine Hopper's, in place of a ledger: one file to a
% notebook, filed under transcripts/notebooks/, the argument the site's slug
% (« battle-of-wash-sq »). Sets the ledger to « notebooks » for everything
% that reads it, so the title block and the watermark behave as they do for a
% batch; the rendering checks the sheets against the notebook's own listing.
\newcommand{\notebook}[1]{\gdef\hopLedger{notebooks}\gdef\hopNotebook{#1}}
\gdef\hopNotebook{}
% The Whitney's accession number for the physical volume — 96.208 … 96.213.
% This is what a citation of the object cites.
\newcommand{\objectnumber}[1]{\gdef\hopObject{#1}}
% The batch this file covers. Leaving it out drops the « (batch N) » from the
% title block rather than printing a number that would be false.
\newcommand{\batch}[1]{\gdef\hopBatch{#1}}
\newcommand{\sheets}[2]{\gdef\hopFirst{#1}\gdef\hopLast{#2}}

% The Whitney's date range for the volume, copied verbatim from its resource
% record — never inferred here, never narrowed to the years this batch
% happens to cover.
\newcommand{\dating}[1]{\gdef\hopDating{#1}}

% The legal notice, printed in the title block and repeated as a diagonal
% watermark on every page of the PDF.
%
% This is not boilerplate and it is not optional. The ledgers are © Heirs of
% Josephine N. Hopper, licensed by the Artists Rights Society; the Whitney
% records the object rights as transferred to the Museum. A transcription
% reproduces Josephine Hopper's text. Every file must therefore say what it is
% on its own face, not only in a README that does not travel with the PDF.
% A file without this line gets no watermark, so omitting it is visible in the
% output rather than silent.
\newcommand{\watermark}[1]{\gdef\hopWatermark{#1}}

\gdef\hopLedger{?}\gdef\hopLedgerTitle{}\gdef\hopObject{}\gdef\hopBatch{}%
\gdef\hopFirst{?}\gdef\hopLast{?}\gdef\hopDating{}\gdef\hopWatermark{}

% --- Keywords ---------------------------------------------------------------
%
% Closing a transcription: the vocabulary under which to search for what these
% leaves record — dealers, exhibiting societies, media, places.
% `npm run manifest` extracts them to tag the whole ledger. This line is the
% single source of the tags; there is no tags file, so no tag can describe
% leaves nobody has read.
%
% It is the one thing in the file that is not on the paper, and it earns its
% place: a reader looking for Keppel needs some way in, and a tag written by
% whoever read the sheets is the only kind that cannot describe unread ones.
\newcommand{\keywords}[1]{\par\smallskip\noindent
  {\footnotesize\textsc{Keywords} --- \textit{#1}}\par}

% --- The anchor -------------------------------------------------------------

% The sheet begins here. Two arguments, and both are needed.
%
% #1 is the Whitney's **resource ref** — 16853 — which is the sheet's whole
% address: it is what the image URL is built from and what turns the facsimile
% as the transcript is scrolled. #2 is the number **written on the paper**, or
% empty where nobody wrote one, which is the case for every cover, flyleaf,
% index and loose insertion — about one sheet in nine.
%
% Keeping both is the whole of the numbering problem in this archive. Three
% sequences overlap: the Hoppers' own leaf numbers, the Whitney's order of
% photographs, and the resource refs, which are in no order at all. Only the
% first is citable and only the third is addressable, so a transcription that
% recorded one of them would be either uncitable or unopenable.
%
% `scripts/render.mjs` checks #2 against what `scripts/catalogue.mjs` parsed
% out of the Whitney's descriptor and fails on a disagreement, because a
% mistyped leaf number silently moves an entry to another year.
\newcommand{\sheet}[2]{%
  \marginnote{\footnotesize\color{gray!70}\textsf{\ifx\relax#2\relax---\else#2\fi}}%
  \label{sheet:#1}\ignorespaces}

% --- The critical apparatus -------------------------------------------------

% Illegible: never guessed. On a ledger leaf this is most often a figure, and
% a guessed figure is a fabricated transaction.
\newcommand{\ill}{\textcolor{red!60!black}{[\,\ldots\,]}}

% A doubtful reading: the word or figure is given, and flagged as uncertain.
%
% Underlining belongs to this macro and to nothing else. Jo Hopper underlines
% her own headings --- « Etchings », « Etching Plates », « Early oil » --- and
% those are set with \emph{} instead, because a reader must be able to tell a
% doubtful reading from her emphasis at a glance and the two cannot both be
% underlines. Which of them keeps the underline is decided by consequence: an
% emphasis shown as italic loses nothing, and a doubtful figure shown as a
% certain one loses everything.
\newcommand{\uncertain}[1]{\uline{#1}}

% An editorial addition — an expanded abbreviation, an implied dollar sign.
\newcommand{\add}[1]{\textcolor{blue!50!black}{[#1]}}

% Struck out in the book. Jo Hopper struck a great deal: a price revised, a
% buyer who withdrew, an exhibition that fell through. What she crossed out is
% part of the record and often the more interesting half of it.
\newcommand{\struck}[1]{\sout{#1}}

% The ink it is written in, where the ink says something.
%
% Book IV is a state machine and the colour is its status field: charges in
% black, receipts in red, the sums in pencil, a rule that holds from leaf 7 to
% the end of the volume. A transcription that records the words and the
% figures and nothing about their colour throws that field away, and the
% accounts then have to guess a receipt from the word « rec'd » --- which the
% volume stops writing on leaf 157. So the ink is recorded, per cell, as what
% is on the paper and not as what it means: \ink{red}{rec'd by check} says
% the words are red, and it is `scripts/accounts.mjs` that says red is money
% received. Black is the default and is not marked. The inks are the three
% the volume uses besides black; a fourth would be added here and in
% `scripts/render.mjs` and `scripts/tei.mjs` in the same commit.
% The file is \input, not \usepackage'd, so @ is not a letter here without saying so.
\makeatletter
\@namedef{hop@ink@red}{red!70!black}
\@namedef{hop@ink@pencil}{gray!55!black}
\@namedef{hop@ink@blue}{blue!55!black}
\newcommand{\ink}[2]{%
  \ifcsname hop@ink@#1\endcsname
    \textcolor{\csname hop@ink@#1\endcsname}{#2}%
  \else
    \PackageError{hopper}{\string\ink{#1}: unknown ink}{The inks are red, pencil and blue.}%
  \fi}
\makeatother

% A rule drawn across the money column above this row's figure.
%
% The writer rules off a run of items and writes their sum under the rule,
% and the rule is the whole of what says « this figure is a sum ». It stands
% at the head of the row the sum is on: \ruledoff & & 4355 & 00. On paper it
% is a line under the last two columns of the four a Book IV leaf has; if
% another volume ever rules a different column, the macro grows an argument.
\newcommand{\ruledoff}{\cline{3-4}}

% A diary entry begins. #1 is the date as she wrote it — « Mar. 12 »,
% « Sunday » — and #2 the date the transcriber assigns it, as a year, a month
% and a day (1947-03-12), or as much of one as the leaf allows. The first is
% what is on the paper; the second is the one editorial claim a diary needs
% that a ledger does not, because a ledger's leaf is its own address and an
% entry's date is its subject. Both are kept, apart, so the claim can be
% argued with. Used only in the notebooks; a ledger has no entries.
\newcommand{\entry}[2]{\par\medskip\noindent
  {\bfseries #1}\ifx\relax#2\relax\else\hspace{0.6em}{\footnotesize\color{gray!60!black}\textsf{#2}}\fi
  \par\nopagebreak\smallskip}

% The transcriber's note. Ours, not theirs.
\newcommand{\note}[1]{\par\smallskip{\small\color{gray!60!black}\textit{[#1]}}\par\smallskip}

% A marginal note **of theirs**, distinct from the body — and distinct from
% \note{} for exactly the reason \note{} exists.
\newcommand{\marginal}[1]{\par\smallskip\noindent\hspace{1em}%
  {\small\color{gray!60!black}#1}\par\smallskip}

% --- What is particular to this archive -------------------------------------

% Whose hand wrote it.
%
% This macro has no counterpart in the Grothendieck preamble, and it is the
% single most important thing in this one. These books were written by two
% people and annotated by more. Edward Hopper drew the work and wrote its
% title and size; Josephine Nivison Hopper wrote everything else — the
% descriptions, the anecdotes about people he refused to discuss, the prices,
% the buyers, the columns; and later hands, curatorial and unidentified, added
% to the leaves after 1967.
%
% A transcription that flattens the three has destroyed the archive's chief
% documentary value, which is that it records *two* working lives and the
% division of labour between them. #1 is `edward`, `jo`, `later` or
% `unidentified`; `unidentified` is the right answer far more often than a
% transcriber's confidence suggests, and it is not a failure to use it.
\newcommand{\hand}[2]{%
  \ifx\relax#1\relax#2\else
    \textcolor{gray!55!black}{\footnotesize\textsf{[#1]}}\,#2%
  \fi}

% Edward Hopper's ink record sketch stands here.
%
% The argument is **what is inscribed on or beside the sketch** — a title, a
% dimension, a note — and never a description of the image. Describing the
% drawing would be writing a new document: the sketch is on the site, one pane
% away, at the resolution the Whitney publishes, and prose about it competes
% with looking at it. Leave the argument empty where nothing is written.
% The mark is drawn rather than set from a symbol font: there is no
% mathematics anywhere in this archive, so no maths package is loaded, and
% $\square$ would be an undefined control sequence — which is exactly how it
% failed the first time the PDFs were compiled.
\newcommand{\sketchmark}{\raisebox{0.15ex}{\color{gray!50!black}%
  \tikz[baseline]{\draw[line width=0.4pt] (0,0) rectangle (1.1ex,1.1ex);}}}
\newcommand{\sketch}[1]{\par\smallskip\noindent
  {\small\color{gray!50!black}\sketchmark~\textsf{ink sketch}%
   \ifx\relax#1\relax\else\ --- #1\fi}\par\smallskip}

% Something printed and pasted to the leaf — a reproduction cut from a
% catalogue, a review, a dealer's label. The argument transcribes its printed
% text. These are the only places in the archive where the words are nobody's
% handwriting, and they date the leaf, which handwriting does not.
\newcommand{\clipping}[1]{\par\smallskip\noindent
  {\small\color{gray!50!black}\textsf{[clipping]}\ #1}\par\smallskip}

% A work's block opens here, under the title **exactly as the leaf gives it** —
% « Les Deux Pigeon » and « Les Poillus » stay misspelled, because the
% misspelling is Hopper's and is how the entry is found.
\newcommand{\work}[1]{\par\medskip\noindent{\bfseries #1}\par\nopagebreak\smallskip}

% --- The ruled table --------------------------------------------------------
%
% Jo Hopper ruled these columns herself and kept them for forty years; the
% table is not a presentation of the content, it *is* the content. So it is
% transcribed as a table and rendered as one, on screen and on paper, from one
% source.
%
% #1 is the column specification; #2 is the header row, `&`-separated like any
% other. A cell she left blank is left blank; a cell that cannot be read takes
% \ill{}.
%
% **Use `Y{...}` for any column that can hold a sentence** — an exhibition, a
% dealer's terms, a note — and `l`, `r` or `c` only for dates, figures and
% short names. See the note on \newcolumntype{Y} above for why that is not a
% matter of taste.
%
% The header row is emitted **bare**, between rules, and is deliberately not
% wrapped in \textbf{}. It cannot be: an `&` inside a braced group is not an
% alignment tab, so `\textbf{Date & Exhibitions}` fails at the first header
% with « Forbidden control sequence found while scanning use of
% \check@nocorr@ » — which is TeX saying, at some distance, that the tabs are
% in the wrong place. Rules carry the header instead, which is how a scholarly
% table sets one anyway. (The reading view has no such constraint and renders
% the same row as <th>.)
\newenvironment{ledgertable}[2]
  {\par\smallskip\small
   \begin{longtable}{#1}
   \hline
   #2 \\
   \hline
   \endhead}
  {\end{longtable}\par\smallskip}

% --- The appendix of notes --------------------------------------------------
%
% Everything above the appendix is on the paper. Nothing in it is.
%
% A note says what a named source says about a work these leaves record --- what
% the holding museum measures the canvas at, which other leaf carries the other
% half of a transaction --- and it is the only prose in the file that a reader
% cannot check against the photograph. On the site that is handled by putting
% the note beside a link to its source; on paper a link is not enough, because
% a PDF outlives the site it was downloaded from and travels away from it. So
% every claim here prints its source in full, name and address both, and every
% note prints the day it was written and what wrote it.
%
% The appendix is assembled by `scripts/pdf.mjs` from
% `src/content/work-notes.json` at compile time and is in no `.tex` file. That
% is the point of it. The transcription is the edition; a note is a later and
% weaker kind of writing *about* the edition, and keeping the two in one source
% file would eventually put a sentence somebody inferred inside the document
% whose whole claim on a reader is that nothing in it was inferred.
%
% Only the notes on works this batch's leaves name are printed, so the appendix
% follows the batch rather than repeating the archive, and a batch that has
% earned no notes yet gets no appendix at all --- which is a true statement
% about how much has been checked, and the same statement `work-notes.json`
% makes by being nearly empty.

% A URL is the one string in this archive that can run a hundred characters
% without a space in it, and an overfull box fails the build here because in a
% transcription an overfull box hides content. So the appendix tells TeX where
% a URL may break and sets its apparatus ragged right, rather than leaving the
% line to overrun in a file nobody recompiles.
\def\UrlBreaks{\do\/\do\-\do\.\do\_\do\?\do\&\do\=\do\+}
\Urlmuskip=0mu plus 1mu

\newcommand{\notesappendix}{\clearpage
  \par\noindent{\large\bfseries Notes on the works}\par\smallskip
  \noindent{\footnotesize\color{gray!60!black}\raggedright
    Not on the paper, and not part of the transcription. Each note below
    concerns a work named on a leaf of this batch and says what a named source
    says about it --- a museum's own catalogue record, or another leaf of these
    ledgers. Every sentence carries the source that states it, printed with its
    address so that it can be followed from a copy of this file alone. Where a
    ledger and a museum disagree the disagreement is printed rather than
    resolved.\par}\par\medskip}

% One work: the title as the leaves give it, then the note.
\newcommand{\worknote}[2]{\par\medskip\noindent{\bfseries #1}\par\nopagebreak
  \smallskip\noindent #2\par}

% One claim under it: what is asserted, and who asserts it.
\newcommand{\workclaim}[2]{\par\smallskip\noindent
  {\footnotesize\color{gray!60!black}\raggedright\hangindent=1.6em\hangafter=0
   \hspace{1em}#1\ --- #2\par}}

% Who wrote the note and when, closing the block. A note is a model's reading
% of retrieved records and is weaker evidence than anything above it; the file
% records that rather than letting the appendix pass for the edition.
\newcommand{\worknoteby}[2]{\par\smallskip\noindent
  {\scriptsize\color{gray!55!black}\hspace{1em}Written #1 by #2.\par}\par\smallskip}

% --- Title ------------------------------------------------------------------

% The watermark is drawn on the shipout background so it sits under the text of
% every page, diagonal and light enough to leave the record readable.
\AddToHook{shipout/background}{%
  \ifx\hopWatermark\empty\else
    \begin{tikzpicture}[remember picture, overlay]
      \node[rotate=52, scale=3.4, align=center, text=gray!18,
            font=\sffamily\bfseries]
        at (current page.center) {\hopWatermark};
    \end{tikzpicture}%
  \fi}

\AtBeginDocument{%
  \begin{center}
    {\large\bfseries Edward and Josephine Hopper --- \hopLedgerTitle}\\[2pt]
    \ifx\hopObject\empty\else
      {\small Whitney Museum of American Art, \hopObject}\\[4pt]
    \fi
    {\small sheets \hopFirst--\hopLast
      \ifx\hopBatch\empty\else\ (batch \hopBatch)\fi}\\[2pt]
    \ifx\hopDating\empty\else
      {\small\color{gray!60!black}The Whitney's dating: \hopDating}\\[2pt]
    \fi
    \ifx\hopWatermark\empty\else
      {\small\bfseries\color{red!45!gray}\hopWatermark}\\[2pt]
    \fi
    {\footnotesize\color{gray!60!black}Transcribed from the digitised sheets published by the
     Whitney Museum of American Art.\\
     \textcopyright{} Heirs of Josephine N. Hopper, licensed by Artists Rights Society (ARS),
     New York.\\
     Uncertain readings are underlined, illegible passages marked [\,\ldots\,],
     editorial additions bracketed.}
  \end{center}
  \vspace{1em}}

\endinput
 rather than \usepackage, and \ProvidesPackage in
% an \input-ed file makes LaTeX warn on every compile that it "requested
% package `'" and got `hopper'. A warning nobody can act on is a warning
% everybody learns to scroll past, which is how the real ones get missed.
\ProvidesFile{hopper.sty}[2026/09/05 Transcriptions of the Hopper artist's ledgers]

\usepackage{array}
% longtable, not tabularx: a leaf of thirty-one exhibition lines must break
% across pages, and tabularx cannot be wrapped in an environment of our own —
% it scans the source for a literal \end{tabularx} and dies on \end{ledgertable}
% with « File ended while scanning use of \TX@get@body ». longtable has no such
% scan, and the wrapping column type below does the job tabularx's X would.
\usepackage{longtable}

% A wrapping column, `Y{0.3}` being three tenths of the measure.
%
% Prose columns must be declared this way and not as `l`. A table of `l`
% columns is set to its natural width, which can be wider than the page, and
% TeX issues **no warning at all** because nothing ever asked the row to fit:
% the first compile of Book I produced a page whose right-hand column ran off
% the paper with a clean log. A fixed-width column wraps, and if it still
% cannot fit, TeX says so — which is what `npm run pdf` then refuses to build
% over.
%
% Leave about 0.03 of the measure per column for the inter-column space: five
% columns can share 0.85, six can share 0.80. Getting it wrong is not a
% guessing game — `npm run pdf` reports the overfull box and how many points
% too wide it is.
\newcolumntype{Y}[1]{>{\raggedright\arraybackslash}p{#1\linewidth}}
\usepackage[english]{babel}
\usepackage{fontspec}
\usepackage{geometry}
\geometry{a4paper,margin=25mm,marginparwidth=28mm}
\usepackage{marginnote}
\usepackage{xcolor}
\usepackage{tikz}
% ulem, not soul. `soul`'s \st and \ul are built on a character-by-character
% scan that XeLaTeX's font machinery breaks: the document fails with an
% undefined control sequence rather than degrading. `normalem` stops it
% hijacking \emph, which the transcriptions use for underlining of Jo Hopper's
% own.
\usepackage[normalem]{ulem}
% ulem sets each underlined word in a box TeX can neither hyphenate nor break
% after, so a paragraph dense in \uncertain{} readings can reach a state where
% no legal breakpoint exists. \emergencystretch lets TeX loosen it on a third
% pass instead of overrunning: an overfull line in a transcription hides
% content.
\setlength{\emergencystretch}{3em}
\usepackage{hyperref}
\hypersetup{colorlinks=true,linkcolor=black,urlcolor=black,pdfborder={0 0 0}}

% --- Batch metadata ---------------------------------------------------------
% Taken as they stand by the rendering script, which shows them at the head of
% the reading view. They fix what the file claims to cover: without them a
% transcription floats unmoored in an archive of five hundred sheets.

\newcommand{\ledger}[1]{\gdef\hopLedger{#1}}
\newcommand{\ledgertitle}[1]{\gdef\hopLedgerTitle{#1}}
% A notebook of Josephine Hopper's, in place of a ledger: one file to a
% notebook, filed under transcripts/notebooks/, the argument the site's slug
% (« battle-of-wash-sq »). Sets the ledger to « notebooks » for everything
% that reads it, so the title block and the watermark behave as they do for a
% batch; the rendering checks the sheets against the notebook's own listing.
\newcommand{\notebook}[1]{\gdef\hopLedger{notebooks}\gdef\hopNotebook{#1}}
\gdef\hopNotebook{}
% The Whitney's accession number for the physical volume — 96.208 … 96.213.
% This is what a citation of the object cites.
\newcommand{\objectnumber}[1]{\gdef\hopObject{#1}}
% The batch this file covers. Leaving it out drops the « (batch N) » from the
% title block rather than printing a number that would be false.
\newcommand{\batch}[1]{\gdef\hopBatch{#1}}
\newcommand{\sheets}[2]{\gdef\hopFirst{#1}\gdef\hopLast{#2}}

% The Whitney's date range for the volume, copied verbatim from its resource
% record — never inferred here, never narrowed to the years this batch
% happens to cover.
\newcommand{\dating}[1]{\gdef\hopDating{#1}}

% The legal notice, printed in the title block and repeated as a diagonal
% watermark on every page of the PDF.
%
% This is not boilerplate and it is not optional. The ledgers are © Heirs of
% Josephine N. Hopper, licensed by the Artists Rights Society; the Whitney
% records the object rights as transferred to the Museum. A transcription
% reproduces Josephine Hopper's text. Every file must therefore say what it is
% on its own face, not only in a README that does not travel with the PDF.
% A file without this line gets no watermark, so omitting it is visible in the
% output rather than silent.
\newcommand{\watermark}[1]{\gdef\hopWatermark{#1}}

\gdef\hopLedger{?}\gdef\hopLedgerTitle{}\gdef\hopObject{}\gdef\hopBatch{}%
\gdef\hopFirst{?}\gdef\hopLast{?}\gdef\hopDating{}\gdef\hopWatermark{}

% --- Keywords ---------------------------------------------------------------
%
% Closing a transcription: the vocabulary under which to search for what these
% leaves record — dealers, exhibiting societies, media, places.
% `npm run manifest` extracts them to tag the whole ledger. This line is the
% single source of the tags; there is no tags file, so no tag can describe
% leaves nobody has read.
%
% It is the one thing in the file that is not on the paper, and it earns its
% place: a reader looking for Keppel needs some way in, and a tag written by
% whoever read the sheets is the only kind that cannot describe unread ones.
\newcommand{\keywords}[1]{\par\smallskip\noindent
  {\footnotesize\textsc{Keywords} --- \textit{#1}}\par}

% --- The anchor -------------------------------------------------------------

% The sheet begins here. Two arguments, and both are needed.
%
% #1 is the Whitney's **resource ref** — 16853 — which is the sheet's whole
% address: it is what the image URL is built from and what turns the facsimile
% as the transcript is scrolled. #2 is the number **written on the paper**, or
% empty where nobody wrote one, which is the case for every cover, flyleaf,
% index and loose insertion — about one sheet in nine.
%
% Keeping both is the whole of the numbering problem in this archive. Three
% sequences overlap: the Hoppers' own leaf numbers, the Whitney's order of
% photographs, and the resource refs, which are in no order at all. Only the
% first is citable and only the third is addressable, so a transcription that
% recorded one of them would be either uncitable or unopenable.
%
% `scripts/render.mjs` checks #2 against what `scripts/catalogue.mjs` parsed
% out of the Whitney's descriptor and fails on a disagreement, because a
% mistyped leaf number silently moves an entry to another year.
\newcommand{\sheet}[2]{%
  \marginnote{\footnotesize\color{gray!70}\textsf{\ifx\relax#2\relax---\else#2\fi}}%
  \label{sheet:#1}\ignorespaces}

% --- The critical apparatus -------------------------------------------------

% Illegible: never guessed. On a ledger leaf this is most often a figure, and
% a guessed figure is a fabricated transaction.
\newcommand{\ill}{\textcolor{red!60!black}{[\,\ldots\,]}}

% A doubtful reading: the word or figure is given, and flagged as uncertain.
%
% Underlining belongs to this macro and to nothing else. Jo Hopper underlines
% her own headings --- « Etchings », « Etching Plates », « Early oil » --- and
% those are set with \emph{} instead, because a reader must be able to tell a
% doubtful reading from her emphasis at a glance and the two cannot both be
% underlines. Which of them keeps the underline is decided by consequence: an
% emphasis shown as italic loses nothing, and a doubtful figure shown as a
% certain one loses everything.
\newcommand{\uncertain}[1]{\uline{#1}}

% An editorial addition — an expanded abbreviation, an implied dollar sign.
\newcommand{\add}[1]{\textcolor{blue!50!black}{[#1]}}

% Struck out in the book. Jo Hopper struck a great deal: a price revised, a
% buyer who withdrew, an exhibition that fell through. What she crossed out is
% part of the record and often the more interesting half of it.
\newcommand{\struck}[1]{\sout{#1}}

% The ink it is written in, where the ink says something.
%
% Book IV is a state machine and the colour is its status field: charges in
% black, receipts in red, the sums in pencil, a rule that holds from leaf 7 to
% the end of the volume. A transcription that records the words and the
% figures and nothing about their colour throws that field away, and the
% accounts then have to guess a receipt from the word « rec'd » --- which the
% volume stops writing on leaf 157. So the ink is recorded, per cell, as what
% is on the paper and not as what it means: \ink{red}{rec'd by check} says
% the words are red, and it is `scripts/accounts.mjs` that says red is money
% received. Black is the default and is not marked. The inks are the three
% the volume uses besides black; a fourth would be added here and in
% `scripts/render.mjs` and `scripts/tei.mjs` in the same commit.
% The file is \input, not \usepackage'd, so @ is not a letter here without saying so.
\makeatletter
\@namedef{hop@ink@red}{red!70!black}
\@namedef{hop@ink@pencil}{gray!55!black}
\@namedef{hop@ink@blue}{blue!55!black}
\newcommand{\ink}[2]{%
  \ifcsname hop@ink@#1\endcsname
    \textcolor{\csname hop@ink@#1\endcsname}{#2}%
  \else
    \PackageError{hopper}{\string\ink{#1}: unknown ink}{The inks are red, pencil and blue.}%
  \fi}
\makeatother

% A rule drawn across the money column above this row's figure.
%
% The writer rules off a run of items and writes their sum under the rule,
% and the rule is the whole of what says « this figure is a sum ». It stands
% at the head of the row the sum is on: \ruledoff & & 4355 & 00. On paper it
% is a line under the last two columns of the four a Book IV leaf has; if
% another volume ever rules a different column, the macro grows an argument.
\newcommand{\ruledoff}{\cline{3-4}}

% A diary entry begins. #1 is the date as she wrote it — « Mar. 12 »,
% « Sunday » — and #2 the date the transcriber assigns it, as a year, a month
% and a day (1947-03-12), or as much of one as the leaf allows. The first is
% what is on the paper; the second is the one editorial claim a diary needs
% that a ledger does not, because a ledger's leaf is its own address and an
% entry's date is its subject. Both are kept, apart, so the claim can be
% argued with. Used only in the notebooks; a ledger has no entries.
\newcommand{\entry}[2]{\par\medskip\noindent
  {\bfseries #1}\ifx\relax#2\relax\else\hspace{0.6em}{\footnotesize\color{gray!60!black}\textsf{#2}}\fi
  \par\nopagebreak\smallskip}

% The transcriber's note. Ours, not theirs.
\newcommand{\note}[1]{\par\smallskip{\small\color{gray!60!black}\textit{[#1]}}\par\smallskip}

% A marginal note **of theirs**, distinct from the body — and distinct from
% \note{} for exactly the reason \note{} exists.
\newcommand{\marginal}[1]{\par\smallskip\noindent\hspace{1em}%
  {\small\color{gray!60!black}#1}\par\smallskip}

% --- What is particular to this archive -------------------------------------

% Whose hand wrote it.
%
% This macro has no counterpart in the Grothendieck preamble, and it is the
% single most important thing in this one. These books were written by two
% people and annotated by more. Edward Hopper drew the work and wrote its
% title and size; Josephine Nivison Hopper wrote everything else — the
% descriptions, the anecdotes about people he refused to discuss, the prices,
% the buyers, the columns; and later hands, curatorial and unidentified, added
% to the leaves after 1967.
%
% A transcription that flattens the three has destroyed the archive's chief
% documentary value, which is that it records *two* working lives and the
% division of labour between them. #1 is `edward`, `jo`, `later` or
% `unidentified`; `unidentified` is the right answer far more often than a
% transcriber's confidence suggests, and it is not a failure to use it.
\newcommand{\hand}[2]{%
  \ifx\relax#1\relax#2\else
    \textcolor{gray!55!black}{\footnotesize\textsf{[#1]}}\,#2%
  \fi}

% Edward Hopper's ink record sketch stands here.
%
% The argument is **what is inscribed on or beside the sketch** — a title, a
% dimension, a note — and never a description of the image. Describing the
% drawing would be writing a new document: the sketch is on the site, one pane
% away, at the resolution the Whitney publishes, and prose about it competes
% with looking at it. Leave the argument empty where nothing is written.
% The mark is drawn rather than set from a symbol font: there is no
% mathematics anywhere in this archive, so no maths package is loaded, and
% $\square$ would be an undefined control sequence — which is exactly how it
% failed the first time the PDFs were compiled.
\newcommand{\sketchmark}{\raisebox{0.15ex}{\color{gray!50!black}%
  \tikz[baseline]{\draw[line width=0.4pt] (0,0) rectangle (1.1ex,1.1ex);}}}
\newcommand{\sketch}[1]{\par\smallskip\noindent
  {\small\color{gray!50!black}\sketchmark~\textsf{ink sketch}%
   \ifx\relax#1\relax\else\ --- #1\fi}\par\smallskip}

% Something printed and pasted to the leaf — a reproduction cut from a
% catalogue, a review, a dealer's label. The argument transcribes its printed
% text. These are the only places in the archive where the words are nobody's
% handwriting, and they date the leaf, which handwriting does not.
\newcommand{\clipping}[1]{\par\smallskip\noindent
  {\small\color{gray!50!black}\textsf{[clipping]}\ #1}\par\smallskip}

% A work's block opens here, under the title **exactly as the leaf gives it** —
% « Les Deux Pigeon » and « Les Poillus » stay misspelled, because the
% misspelling is Hopper's and is how the entry is found.
\newcommand{\work}[1]{\par\medskip\noindent{\bfseries #1}\par\nopagebreak\smallskip}

% --- The ruled table --------------------------------------------------------
%
% Jo Hopper ruled these columns herself and kept them for forty years; the
% table is not a presentation of the content, it *is* the content. So it is
% transcribed as a table and rendered as one, on screen and on paper, from one
% source.
%
% #1 is the column specification; #2 is the header row, `&`-separated like any
% other. A cell she left blank is left blank; a cell that cannot be read takes
% \ill{}.
%
% **Use `Y{...}` for any column that can hold a sentence** — an exhibition, a
% dealer's terms, a note — and `l`, `r` or `c` only for dates, figures and
% short names. See the note on \newcolumntype{Y} above for why that is not a
% matter of taste.
%
% The header row is emitted **bare**, between rules, and is deliberately not
% wrapped in \textbf{}. It cannot be: an `&` inside a braced group is not an
% alignment tab, so `\textbf{Date & Exhibitions}` fails at the first header
% with « Forbidden control sequence found while scanning use of
% \check@nocorr@ » — which is TeX saying, at some distance, that the tabs are
% in the wrong place. Rules carry the header instead, which is how a scholarly
% table sets one anyway. (The reading view has no such constraint and renders
% the same row as <th>.)
\newenvironment{ledgertable}[2]
  {\par\smallskip\small
   \begin{longtable}{#1}
   \hline
   #2 \\
   \hline
   \endhead}
  {\end{longtable}\par\smallskip}

% --- The appendix of notes --------------------------------------------------
%
% Everything above the appendix is on the paper. Nothing in it is.
%
% A note says what a named source says about a work these leaves record --- what
% the holding museum measures the canvas at, which other leaf carries the other
% half of a transaction --- and it is the only prose in the file that a reader
% cannot check against the photograph. On the site that is handled by putting
% the note beside a link to its source; on paper a link is not enough, because
% a PDF outlives the site it was downloaded from and travels away from it. So
% every claim here prints its source in full, name and address both, and every
% note prints the day it was written and what wrote it.
%
% The appendix is assembled by `scripts/pdf.mjs` from
% `src/content/work-notes.json` at compile time and is in no `.tex` file. That
% is the point of it. The transcription is the edition; a note is a later and
% weaker kind of writing *about* the edition, and keeping the two in one source
% file would eventually put a sentence somebody inferred inside the document
% whose whole claim on a reader is that nothing in it was inferred.
%
% Only the notes on works this batch's leaves name are printed, so the appendix
% follows the batch rather than repeating the archive, and a batch that has
% earned no notes yet gets no appendix at all --- which is a true statement
% about how much has been checked, and the same statement `work-notes.json`
% makes by being nearly empty.

% A URL is the one string in this archive that can run a hundred characters
% without a space in it, and an overfull box fails the build here because in a
% transcription an overfull box hides content. So the appendix tells TeX where
% a URL may break and sets its apparatus ragged right, rather than leaving the
% line to overrun in a file nobody recompiles.
\def\UrlBreaks{\do\/\do\-\do\.\do\_\do\?\do\&\do\=\do\+}
\Urlmuskip=0mu plus 1mu

\newcommand{\notesappendix}{\clearpage
  \par\noindent{\large\bfseries Notes on the works}\par\smallskip
  \noindent{\footnotesize\color{gray!60!black}\raggedright
    Not on the paper, and not part of the transcription. Each note below
    concerns a work named on a leaf of this batch and says what a named source
    says about it --- a museum's own catalogue record, or another leaf of these
    ledgers. Every sentence carries the source that states it, printed with its
    address so that it can be followed from a copy of this file alone. Where a
    ledger and a museum disagree the disagreement is printed rather than
    resolved.\par}\par\medskip}

% One work: the title as the leaves give it, then the note.
\newcommand{\worknote}[2]{\par\medskip\noindent{\bfseries #1}\par\nopagebreak
  \smallskip\noindent #2\par}

% One claim under it: what is asserted, and who asserts it.
\newcommand{\workclaim}[2]{\par\smallskip\noindent
  {\footnotesize\color{gray!60!black}\raggedright\hangindent=1.6em\hangafter=0
   \hspace{1em}#1\ --- #2\par}}

% Who wrote the note and when, closing the block. A note is a model's reading
% of retrieved records and is weaker evidence than anything above it; the file
% records that rather than letting the appendix pass for the edition.
\newcommand{\worknoteby}[2]{\par\smallskip\noindent
  {\scriptsize\color{gray!55!black}\hspace{1em}Written #1 by #2.\par}\par\smallskip}

% --- Title ------------------------------------------------------------------

% The watermark is drawn on the shipout background so it sits under the text of
% every page, diagonal and light enough to leave the record readable.
\AddToHook{shipout/background}{%
  \ifx\hopWatermark\empty\else
    \begin{tikzpicture}[remember picture, overlay]
      \node[rotate=52, scale=3.4, align=center, text=gray!18,
            font=\sffamily\bfseries]
        at (current page.center) {\hopWatermark};
    \end{tikzpicture}%
  \fi}

\AtBeginDocument{%
  \begin{center}
    {\large\bfseries Edward and Josephine Hopper --- \hopLedgerTitle}\\[2pt]
    \ifx\hopObject\empty\else
      {\small Whitney Museum of American Art, \hopObject}\\[4pt]
    \fi
    {\small sheets \hopFirst--\hopLast
      \ifx\hopBatch\empty\else\ (batch \hopBatch)\fi}\\[2pt]
    \ifx\hopDating\empty\else
      {\small\color{gray!60!black}The Whitney's dating: \hopDating}\\[2pt]
    \fi
    \ifx\hopWatermark\empty\else
      {\small\bfseries\color{red!45!gray}\hopWatermark}\\[2pt]
    \fi
    {\footnotesize\color{gray!60!black}Transcribed from the digitised sheets published by the
     Whitney Museum of American Art.\\
     \textcopyright{} Heirs of Josephine N. Hopper, licensed by Artists Rights Society (ARS),
     New York.\\
     Uncertain readings are underlined, illegible passages marked [\,\ldots\,],
     editorial additions bracketed.}
  \end{center}
  \vspace{1em}}

\endinput
 rather than \usepackage, and \ProvidesPackage in
% an \input-ed file makes LaTeX warn on every compile that it "requested
% package `'" and got `hopper'. A warning nobody can act on is a warning
% everybody learns to scroll past, which is how the real ones get missed.
\ProvidesFile{hopper.sty}[2026/09/05 Transcriptions of the Hopper artist's ledgers]

\usepackage{array}
% longtable, not tabularx: a leaf of thirty-one exhibition lines must break
% across pages, and tabularx cannot be wrapped in an environment of our own —
% it scans the source for a literal \end{tabularx} and dies on \end{ledgertable}
% with « File ended while scanning use of \TX@get@body ». longtable has no such
% scan, and the wrapping column type below does the job tabularx's X would.
\usepackage{longtable}

% A wrapping column, `Y{0.3}` being three tenths of the measure.
%
% Prose columns must be declared this way and not as `l`. A table of `l`
% columns is set to its natural width, which can be wider than the page, and
% TeX issues **no warning at all** because nothing ever asked the row to fit:
% the first compile of Book I produced a page whose right-hand column ran off
% the paper with a clean log. A fixed-width column wraps, and if it still
% cannot fit, TeX says so — which is what `npm run pdf` then refuses to build
% over.
%
% Leave about 0.03 of the measure per column for the inter-column space: five
% columns can share 0.85, six can share 0.80. Getting it wrong is not a
% guessing game — `npm run pdf` reports the overfull box and how many points
% too wide it is.
\newcolumntype{Y}[1]{>{\raggedright\arraybackslash}p{#1\linewidth}}
\usepackage[english]{babel}
\usepackage{fontspec}
\usepackage{geometry}
\geometry{a4paper,margin=25mm,marginparwidth=28mm}
\usepackage{marginnote}
\usepackage{xcolor}
\usepackage{tikz}
% ulem, not soul. `soul`'s \st and \ul are built on a character-by-character
% scan that XeLaTeX's font machinery breaks: the document fails with an
% undefined control sequence rather than degrading. `normalem` stops it
% hijacking \emph, which the transcriptions use for underlining of Jo Hopper's
% own.
\usepackage[normalem]{ulem}
% ulem sets each underlined word in a box TeX can neither hyphenate nor break
% after, so a paragraph dense in \uncertain{} readings can reach a state where
% no legal breakpoint exists. \emergencystretch lets TeX loosen it on a third
% pass instead of overrunning: an overfull line in a transcription hides
% content.
\setlength{\emergencystretch}{3em}
\usepackage{hyperref}
\hypersetup{colorlinks=true,linkcolor=black,urlcolor=black,pdfborder={0 0 0}}

% --- Batch metadata ---------------------------------------------------------
% Taken as they stand by the rendering script, which shows them at the head of
% the reading view. They fix what the file claims to cover: without them a
% transcription floats unmoored in an archive of five hundred sheets.

\newcommand{\ledger}[1]{\gdef\hopLedger{#1}}
\newcommand{\ledgertitle}[1]{\gdef\hopLedgerTitle{#1}}
% A notebook of Josephine Hopper's, in place of a ledger: one file to a
% notebook, filed under transcripts/notebooks/, the argument the site's slug
% (« battle-of-wash-sq »). Sets the ledger to « notebooks » for everything
% that reads it, so the title block and the watermark behave as they do for a
% batch; the rendering checks the sheets against the notebook's own listing.
\newcommand{\notebook}[1]{\gdef\hopLedger{notebooks}\gdef\hopNotebook{#1}}
\gdef\hopNotebook{}
% The Whitney's accession number for the physical volume — 96.208 … 96.213.
% This is what a citation of the object cites.
\newcommand{\objectnumber}[1]{\gdef\hopObject{#1}}
% The batch this file covers. Leaving it out drops the « (batch N) » from the
% title block rather than printing a number that would be false.
\newcommand{\batch}[1]{\gdef\hopBatch{#1}}
\newcommand{\sheets}[2]{\gdef\hopFirst{#1}\gdef\hopLast{#2}}

% The Whitney's date range for the volume, copied verbatim from its resource
% record — never inferred here, never narrowed to the years this batch
% happens to cover.
\newcommand{\dating}[1]{\gdef\hopDating{#1}}

% The legal notice, printed in the title block and repeated as a diagonal
% watermark on every page of the PDF.
%
% This is not boilerplate and it is not optional. The ledgers are © Heirs of
% Josephine N. Hopper, licensed by the Artists Rights Society; the Whitney
% records the object rights as transferred to the Museum. A transcription
% reproduces Josephine Hopper's text. Every file must therefore say what it is
% on its own face, not only in a README that does not travel with the PDF.
% A file without this line gets no watermark, so omitting it is visible in the
% output rather than silent.
\newcommand{\watermark}[1]{\gdef\hopWatermark{#1}}

\gdef\hopLedger{?}\gdef\hopLedgerTitle{}\gdef\hopObject{}\gdef\hopBatch{}%
\gdef\hopFirst{?}\gdef\hopLast{?}\gdef\hopDating{}\gdef\hopWatermark{}

% --- Keywords ---------------------------------------------------------------
%
% Closing a transcription: the vocabulary under which to search for what these
% leaves record — dealers, exhibiting societies, media, places.
% `npm run manifest` extracts them to tag the whole ledger. This line is the
% single source of the tags; there is no tags file, so no tag can describe
% leaves nobody has read.
%
% It is the one thing in the file that is not on the paper, and it earns its
% place: a reader looking for Keppel needs some way in, and a tag written by
% whoever read the sheets is the only kind that cannot describe unread ones.
\newcommand{\keywords}[1]{\par\smallskip\noindent
  {\footnotesize\textsc{Keywords} --- \textit{#1}}\par}

% --- The anchor -------------------------------------------------------------

% The sheet begins here. Two arguments, and both are needed.
%
% #1 is the Whitney's **resource ref** — 16853 — which is the sheet's whole
% address: it is what the image URL is built from and what turns the facsimile
% as the transcript is scrolled. #2 is the number **written on the paper**, or
% empty where nobody wrote one, which is the case for every cover, flyleaf,
% index and loose insertion — about one sheet in nine.
%
% Keeping both is the whole of the numbering problem in this archive. Three
% sequences overlap: the Hoppers' own leaf numbers, the Whitney's order of
% photographs, and the resource refs, which are in no order at all. Only the
% first is citable and only the third is addressable, so a transcription that
% recorded one of them would be either uncitable or unopenable.
%
% `scripts/render.mjs` checks #2 against what `scripts/catalogue.mjs` parsed
% out of the Whitney's descriptor and fails on a disagreement, because a
% mistyped leaf number silently moves an entry to another year.
\newcommand{\sheet}[2]{%
  \marginnote{\footnotesize\color{gray!70}\textsf{\ifx\relax#2\relax---\else#2\fi}}%
  \label{sheet:#1}\ignorespaces}

% --- The critical apparatus -------------------------------------------------

% Illegible: never guessed. On a ledger leaf this is most often a figure, and
% a guessed figure is a fabricated transaction.
\newcommand{\ill}{\textcolor{red!60!black}{[\,\ldots\,]}}

% A doubtful reading: the word or figure is given, and flagged as uncertain.
%
% Underlining belongs to this macro and to nothing else. Jo Hopper underlines
% her own headings --- « Etchings », « Etching Plates », « Early oil » --- and
% those are set with \emph{} instead, because a reader must be able to tell a
% doubtful reading from her emphasis at a glance and the two cannot both be
% underlines. Which of them keeps the underline is decided by consequence: an
% emphasis shown as italic loses nothing, and a doubtful figure shown as a
% certain one loses everything.
\newcommand{\uncertain}[1]{\uline{#1}}

% An editorial addition — an expanded abbreviation, an implied dollar sign.
\newcommand{\add}[1]{\textcolor{blue!50!black}{[#1]}}

% Struck out in the book. Jo Hopper struck a great deal: a price revised, a
% buyer who withdrew, an exhibition that fell through. What she crossed out is
% part of the record and often the more interesting half of it.
\newcommand{\struck}[1]{\sout{#1}}

% The ink it is written in, where the ink says something.
%
% Book IV is a state machine and the colour is its status field: charges in
% black, receipts in red, the sums in pencil, a rule that holds from leaf 7 to
% the end of the volume. A transcription that records the words and the
% figures and nothing about their colour throws that field away, and the
% accounts then have to guess a receipt from the word « rec'd » --- which the
% volume stops writing on leaf 157. So the ink is recorded, per cell, as what
% is on the paper and not as what it means: \ink{red}{rec'd by check} says
% the words are red, and it is `scripts/accounts.mjs` that says red is money
% received. Black is the default and is not marked. The inks are the three
% the volume uses besides black; a fourth would be added here and in
% `scripts/render.mjs` and `scripts/tei.mjs` in the same commit.
% The file is \input, not \usepackage'd, so @ is not a letter here without saying so.
\makeatletter
\@namedef{hop@ink@red}{red!70!black}
\@namedef{hop@ink@pencil}{gray!55!black}
\@namedef{hop@ink@blue}{blue!55!black}
\newcommand{\ink}[2]{%
  \ifcsname hop@ink@#1\endcsname
    \textcolor{\csname hop@ink@#1\endcsname}{#2}%
  \else
    \PackageError{hopper}{\string\ink{#1}: unknown ink}{The inks are red, pencil and blue.}%
  \fi}
\makeatother

% A rule drawn across the money column above this row's figure.
%
% The writer rules off a run of items and writes their sum under the rule,
% and the rule is the whole of what says « this figure is a sum ». It stands
% at the head of the row the sum is on: \ruledoff & & 4355 & 00. On paper it
% is a line under the last two columns of the four a Book IV leaf has; if
% another volume ever rules a different column, the macro grows an argument.
\newcommand{\ruledoff}{\cline{3-4}}

% A diary entry begins. #1 is the date as she wrote it — « Mar. 12 »,
% « Sunday » — and #2 the date the transcriber assigns it, as a year, a month
% and a day (1947-03-12), or as much of one as the leaf allows. The first is
% what is on the paper; the second is the one editorial claim a diary needs
% that a ledger does not, because a ledger's leaf is its own address and an
% entry's date is its subject. Both are kept, apart, so the claim can be
% argued with. Used only in the notebooks; a ledger has no entries.
\newcommand{\entry}[2]{\par\medskip\noindent
  {\bfseries #1}\ifx\relax#2\relax\else\hspace{0.6em}{\footnotesize\color{gray!60!black}\textsf{#2}}\fi
  \par\nopagebreak\smallskip}

% The transcriber's note. Ours, not theirs.
\newcommand{\note}[1]{\par\smallskip{\small\color{gray!60!black}\textit{[#1]}}\par\smallskip}

% A marginal note **of theirs**, distinct from the body — and distinct from
% \note{} for exactly the reason \note{} exists.
\newcommand{\marginal}[1]{\par\smallskip\noindent\hspace{1em}%
  {\small\color{gray!60!black}#1}\par\smallskip}

% --- What is particular to this archive -------------------------------------

% Whose hand wrote it.
%
% This macro has no counterpart in the Grothendieck preamble, and it is the
% single most important thing in this one. These books were written by two
% people and annotated by more. Edward Hopper drew the work and wrote its
% title and size; Josephine Nivison Hopper wrote everything else — the
% descriptions, the anecdotes about people he refused to discuss, the prices,
% the buyers, the columns; and later hands, curatorial and unidentified, added
% to the leaves after 1967.
%
% A transcription that flattens the three has destroyed the archive's chief
% documentary value, which is that it records *two* working lives and the
% division of labour between them. #1 is `edward`, `jo`, `later` or
% `unidentified`; `unidentified` is the right answer far more often than a
% transcriber's confidence suggests, and it is not a failure to use it.
\newcommand{\hand}[2]{%
  \ifx\relax#1\relax#2\else
    \textcolor{gray!55!black}{\footnotesize\textsf{[#1]}}\,#2%
  \fi}

% Edward Hopper's ink record sketch stands here.
%
% The argument is **what is inscribed on or beside the sketch** — a title, a
% dimension, a note — and never a description of the image. Describing the
% drawing would be writing a new document: the sketch is on the site, one pane
% away, at the resolution the Whitney publishes, and prose about it competes
% with looking at it. Leave the argument empty where nothing is written.
% The mark is drawn rather than set from a symbol font: there is no
% mathematics anywhere in this archive, so no maths package is loaded, and
% $\square$ would be an undefined control sequence — which is exactly how it
% failed the first time the PDFs were compiled.
\newcommand{\sketchmark}{\raisebox{0.15ex}{\color{gray!50!black}%
  \tikz[baseline]{\draw[line width=0.4pt] (0,0) rectangle (1.1ex,1.1ex);}}}
\newcommand{\sketch}[1]{\par\smallskip\noindent
  {\small\color{gray!50!black}\sketchmark~\textsf{ink sketch}%
   \ifx\relax#1\relax\else\ --- #1\fi}\par\smallskip}

% Something printed and pasted to the leaf — a reproduction cut from a
% catalogue, a review, a dealer's label. The argument transcribes its printed
% text. These are the only places in the archive where the words are nobody's
% handwriting, and they date the leaf, which handwriting does not.
\newcommand{\clipping}[1]{\par\smallskip\noindent
  {\small\color{gray!50!black}\textsf{[clipping]}\ #1}\par\smallskip}

% A work's block opens here, under the title **exactly as the leaf gives it** —
% « Les Deux Pigeon » and « Les Poillus » stay misspelled, because the
% misspelling is Hopper's and is how the entry is found.
\newcommand{\work}[1]{\par\medskip\noindent{\bfseries #1}\par\nopagebreak\smallskip}

% --- The ruled table --------------------------------------------------------
%
% Jo Hopper ruled these columns herself and kept them for forty years; the
% table is not a presentation of the content, it *is* the content. So it is
% transcribed as a table and rendered as one, on screen and on paper, from one
% source.
%
% #1 is the column specification; #2 is the header row, `&`-separated like any
% other. A cell she left blank is left blank; a cell that cannot be read takes
% \ill{}.
%
% **Use `Y{...}` for any column that can hold a sentence** — an exhibition, a
% dealer's terms, a note — and `l`, `r` or `c` only for dates, figures and
% short names. See the note on \newcolumntype{Y} above for why that is not a
% matter of taste.
%
% The header row is emitted **bare**, between rules, and is deliberately not
% wrapped in \textbf{}. It cannot be: an `&` inside a braced group is not an
% alignment tab, so `\textbf{Date & Exhibitions}` fails at the first header
% with « Forbidden control sequence found while scanning use of
% \check@nocorr@ » — which is TeX saying, at some distance, that the tabs are
% in the wrong place. Rules carry the header instead, which is how a scholarly
% table sets one anyway. (The reading view has no such constraint and renders
% the same row as <th>.)
\newenvironment{ledgertable}[2]
  {\par\smallskip\small
   \begin{longtable}{#1}
   \hline
   #2 \\
   \hline
   \endhead}
  {\end{longtable}\par\smallskip}

% --- The appendix of notes --------------------------------------------------
%
% Everything above the appendix is on the paper. Nothing in it is.
%
% A note says what a named source says about a work these leaves record --- what
% the holding museum measures the canvas at, which other leaf carries the other
% half of a transaction --- and it is the only prose in the file that a reader
% cannot check against the photograph. On the site that is handled by putting
% the note beside a link to its source; on paper a link is not enough, because
% a PDF outlives the site it was downloaded from and travels away from it. So
% every claim here prints its source in full, name and address both, and every
% note prints the day it was written and what wrote it.
%
% The appendix is assembled by `scripts/pdf.mjs` from
% `src/content/work-notes.json` at compile time and is in no `.tex` file. That
% is the point of it. The transcription is the edition; a note is a later and
% weaker kind of writing *about* the edition, and keeping the two in one source
% file would eventually put a sentence somebody inferred inside the document
% whose whole claim on a reader is that nothing in it was inferred.
%
% Only the notes on works this batch's leaves name are printed, so the appendix
% follows the batch rather than repeating the archive, and a batch that has
% earned no notes yet gets no appendix at all --- which is a true statement
% about how much has been checked, and the same statement `work-notes.json`
% makes by being nearly empty.

% A URL is the one string in this archive that can run a hundred characters
% without a space in it, and an overfull box fails the build here because in a
% transcription an overfull box hides content. So the appendix tells TeX where
% a URL may break and sets its apparatus ragged right, rather than leaving the
% line to overrun in a file nobody recompiles.
\def\UrlBreaks{\do\/\do\-\do\.\do\_\do\?\do\&\do\=\do\+}
\Urlmuskip=0mu plus 1mu

\newcommand{\notesappendix}{\clearpage
  \par\noindent{\large\bfseries Notes on the works}\par\smallskip
  \noindent{\footnotesize\color{gray!60!black}\raggedright
    Not on the paper, and not part of the transcription. Each note below
    concerns a work named on a leaf of this batch and says what a named source
    says about it --- a museum's own catalogue record, or another leaf of these
    ledgers. Every sentence carries the source that states it, printed with its
    address so that it can be followed from a copy of this file alone. Where a
    ledger and a museum disagree the disagreement is printed rather than
    resolved.\par}\par\medskip}

% One work: the title as the leaves give it, then the note.
\newcommand{\worknote}[2]{\par\medskip\noindent{\bfseries #1}\par\nopagebreak
  \smallskip\noindent #2\par}

% One claim under it: what is asserted, and who asserts it.
\newcommand{\workclaim}[2]{\par\smallskip\noindent
  {\footnotesize\color{gray!60!black}\raggedright\hangindent=1.6em\hangafter=0
   \hspace{1em}#1\ --- #2\par}}

% Who wrote the note and when, closing the block. A note is a model's reading
% of retrieved records and is weaker evidence than anything above it; the file
% records that rather than letting the appendix pass for the edition.
\newcommand{\worknoteby}[2]{\par\smallskip\noindent
  {\scriptsize\color{gray!55!black}\hspace{1em}Written #1 by #2.\par}\par\smallskip}

% --- Title ------------------------------------------------------------------

% The watermark is drawn on the shipout background so it sits under the text of
% every page, diagonal and light enough to leave the record readable.
\AddToHook{shipout/background}{%
  \ifx\hopWatermark\empty\else
    \begin{tikzpicture}[remember picture, overlay]
      \node[rotate=52, scale=3.4, align=center, text=gray!18,
            font=\sffamily\bfseries]
        at (current page.center) {\hopWatermark};
    \end{tikzpicture}%
  \fi}

\AtBeginDocument{%
  \begin{center}
    {\large\bfseries Edward and Josephine Hopper --- \hopLedgerTitle}\\[2pt]
    \ifx\hopObject\empty\else
      {\small Whitney Museum of American Art, \hopObject}\\[4pt]
    \fi
    {\small sheets \hopFirst--\hopLast
      \ifx\hopBatch\empty\else\ (batch \hopBatch)\fi}\\[2pt]
    \ifx\hopDating\empty\else
      {\small\color{gray!60!black}The Whitney's dating: \hopDating}\\[2pt]
    \fi
    \ifx\hopWatermark\empty\else
      {\small\bfseries\color{red!45!gray}\hopWatermark}\\[2pt]
    \fi
    {\footnotesize\color{gray!60!black}Transcribed from the digitised sheets published by the
     Whitney Museum of American Art.\\
     \textcopyright{} Heirs of Josephine N. Hopper, licensed by Artists Rights Society (ARS),
     New York.\\
     Uncertain readings are underlined, illegible passages marked [\,\ldots\,],
     editorial additions bracketed.}
  \end{center}
  \vspace{1em}}

\endinput


\ledger{dealers}
\ledgertitle{Artist's ledger --- Dealers/Etchings}
\objectnumber{96.213}
\batch{4}
\sheets{37}{48}
\dating{1921--1951}
\watermark{Demonstration edition}

\begin{document}

\section{Leaf 88: Mrs. Albert Sterner, E. \& A. Milch, the American Society of
Graphic Arts}

\sheet{17130}{88}

\hand{jo}{Mrs. Albert Sterner \quad 22 W. 49}

\begin{ledgertable}{Y{0.14}Y{0.06}Y{0.25}Y{0.13}Y{0.08}Y{0.14}}{ & & & Sold & Amt. & Returned}
May 16, 22 & 18 & Train + Bathers & & & Returned \\
 & 18 & Les Deux Pigeons & Ap. 16, 23 & 12.60 & * \\
 & 18 & American Landscape & & & Returned \\
 & 18 & Bull Fight & & & '' \\
 & 18 & Evening Wind & Jan. 12, 23 & 12.60 & \\
 & 18 & Night Shadows & & & Ret. \\
 & 15 & Night in the Park & Jan. 18, 24 & 12.60 & \\
 & 18 & Railroad & & & Ret. \\
 & 18 & East Side Interior & Jan. 12, 23 & 12.60 & \\
 & 18 & Mohegan Boat & & & Ret. \\
Jan. 4, 23 & 18 & East Side Interior & & & Ret. \\
 & 18 & Evening Wind & Ap. 16, 23 & 12.60 & \\
 & & '' \quad '' & & & Ret. \\
Ap. 25, 23. & 22 & Evening Wind & Jan. 18, 24 & 15.40 & \\
'' & 22 & '' \quad '' & Jan. 18, 24 & 15.40 & \\
Nov. 8, 23 & 18 & Aux Fortifications & & & Ret. \\
Jan. 24, 24 & 22 & Evening Wind & Dec. 31, 25 & 14.67 & \\
 & & '' \quad '' & & & Ret. \\
 & 25 & East Side Interior & & & \ill{} \\
 & & '' \quad '' & & & Mar. 16, 26. \\
 & 18 & Mohegan Boat & & & Mar. 16, 26 \\
 & & '' \quad '' & & & '' \\
 & 18 & Night in the Park & & & '' \\
\end{ledgertable}

\note{« Sold \quad Amt. \quad Returned » is the only heading written on the
leaf; the date, price and title columns are unheaded and are left so here. The
Returned column alternates « Returned » and « Ret. » as written. The asterisk
against Les Deux Pigeons is hers, drawn as a six-pointed star, and nothing on
this leaf answers it. The price against Night in the Park is 15 where every
other line in the block is 18 or 22; the second figure is formed differently
from the eights above it and is read as a five. Against the last five rows only
a date stands in the Returned column, with no word: the leaf gives the date and
not what happened on it. The cell marked illegible is a mark rather than a
word, struck through and unreadable.}

\hand{jo}{E. \& A. Milch, Inc. \quad 108 W. 57''}

\note{the Milch account is written in two halves side by side, ruled off from
one another by a vertical stroke, and the rows run in parallel rather than in
sequence. The left half is given first and the right half second.}

\begin{ledgertable}{Y{0.15}Y{0.07}Y{0.27}Y{0.14}Y{0.17}}{ & & & & }
Jan. 25, 22 & 18 & Deux Pigeons & Feb. 5, 38 & Returned \\
 & 18 & Night - the L Train & Oct. 15, 29 & \$12. \\
 & \struck{15} 18 & House Tops. & Oct. 15, 29 & 12. \\
\end{ledgertable}

\begin{ledgertable}{Y{0.07}Y{0.30}Y{0.16}Y{0.17}}{ & & & }
18 & Summer Twilight & & \\
\struck{25} 30 & Evening Wind & Oct. 15, 29 & \$20. \\
\struck{25} 30 & East Side Interior & Oct. 15, 29 & 20. \\
\end{ledgertable}

\hand{jo}{Milch Continued}

\note{« Milch Continued » is written inside a shape drawn round it in the same
ink, on the ruling above the four rows that follow; the rows below it are boxed
off by hand-drawn lines on all four sides. It is the same account resumed seven
years later.}

\begin{ledgertable}{Y{0.15}Y{0.07}Y{0.28}Y{0.28}}{ & & & }
Nov. 27, 29. & 30 & Evening Wind & Feb. 5, 38 \quad Returned \\
'' & 30 & Evening Wind & July 1, 49 \quad '' \\
'' & 30 & Night Shadows & Feb. 5, 38 \quad Return\add{ed} \\
'' & 30 & Night in the Park & Returned July 1, '49 \\
\end{ledgertable}

\begin{ledgertable}{Y{0.15}Y{0.07}Y{0.28}Y{0.28}}{ & & & }
Nov. 27, 29 & 30 & East Side Int. & Returned July 1, '49 \\
'' & 25 & R. R. & '' \quad '' \quad '' \\
'' & 30 & Cat Boat & '' \quad '' \quad '' \\
'' & 30 & Locomotive & '' \quad '' \quad '' \\
'' & 30 & Amer. Landscape & Feb. 5, 38 \quad Return\add{ed} \\
\end{ledgertable}

\marginal{Milch}

\note{« Milch » is written up the leaf, sideways, in the margin to the left of
these five rows, with a brace gathering them. The three ditto marks in the
right-hand column stand under « Returned July 1, '49 ».}

\hand{jo}{Amer. Soc. Graphic Arts}

\begin{ledgertable}{Y{0.14}Y{0.07}Y{0.29}Y{0.28}}{ & & & }
Oct. 12 & \$25 & East Side Interior & Returned \\
 & 18 & Night in the Park & \\
Dec. 8 '' & 25 & East Side Interior & \\
 & 18 & Night in the Park. & \\
\end{ledgertable}

\note{a single brace gathers the four rows and « Returned » is written once
against it, not against the first row alone. The year is not given for either
date. A red pencil stroke crosses this block diagonally.}

\section{Leaf 89: Sidney Phillips, Babcock, the Print Rooms, Carl Smalley,
Brown Robertson, Robertson \uncertain{Deschamps}}

\sheet{16668}{89}

\hand{jo}{\emph{Sidney Phillips} \quad 7 W. 51''}

\begin{ledgertable}{Y{0.16}Y{0.32}Y{0.15}Y{0.08}Y{0.09}}{ & & Sold & amt. & }
\struck{Mar. 3, 23.} & Evening Wind & Mar. 3, 23 & 9 & \\
 & Night in the Park & '' & 9 & \\
 & Railroad & & 9 & * \\
 & East Side Interior & & 9 & ** \\
 & '' \quad '' \quad '' & & 9 & ** \\
 & '' \quad '' \quad '' & Ap. 3, 23 & 9 & \\
 & Night Shadows & June 5, 23 & 9 & \\
 & Night in Park - sold \& \add{not} paid for. & & & * \\
\end{ledgertable}

\note{« Sold \quad amt. » are the only headings. The date at the head of the
first row is struck through in the same ink; the same date stands unstruck in
the Sold column beside it. One asterisk and two double asterisks stand in the
last column and nothing on the leaf answers them. « Night in Park - sold \&
not paid for. » is written across the ruling of the last row, and the price
and amount columns are empty against it. A red pencil stroke runs from this
heading down through the leaf.}

\hand{jo}{\emph{Babcock Art Galleries} \quad 19 E. 49''}

\note{the Babcock account is written in two halves side by side, ruled off by
vertical strokes, and the halves are not in sequence with one another. The left
half is given first.}

\begin{ledgertable}{Y{0.16}Y{0.10}Y{0.28}Y{0.08}Y{0.18}}{ & & & & }
Ap. 3, '31 & 30 \struck{25} & East Side Interior & 20 & Feb. 6, '30 \\
 & 30 \struck{25} & Evening Wind & 20 & Mar. 10, '30. \\
Oct. 30, 30 & 30 & Night in the Park & \struck{20 \quad Sept. 20, '38} & \\
\end{ledgertable}

\begin{ledgertable}{Y{0.15}Y{0.07}Y{0.26}Y{0.14}Y{0.15}}{ & & & & }
Oct. 30, 30 & \$30 & Cat Boat & & \\
'' & '' & Amer. Landscape & & \\
'' & '' & Evening Wind & Oct. 22, 43 & \$15 \\
'' & '' & Night Shadows & Sept. 30, 48 & \$20 \\
\end{ledgertable}

\marginal{Sold for \ill{} by mistake}

\note{« Sold for » and a struck figure, with « by mistake » below, are written
small in the gap between the price and title columns beside the Evening Wind
row. The figure is struck through and cannot be read. In the left half, the
price on the first two rows is a 30 written above a 25 that has been struck;
the whole of the Night in the Park cell - the 20 and the date after it - is
struck through in a single heavy stroke.}

\hand{jo}{\emph{Print Rooms} - 540 Sutter St., San Francisco}

\begin{ledgertable}{Y{0.08}Y{0.32}Y{0.16}Y{0.10}}{ & & & }
18 & Night in the Park & Mar. 19, 24 & 9. \\
\end{ledgertable}

\hand{jo}{\emph{Carl Smalley}, MacPherson, Kansas \quad Dead?}

\begin{ledgertable}{Y{0.40}Y{0.12}Y{0.08}Y{0.08}}{ & & & }
Les Poilus (somewhere in France) & 1919. & 8 & * \\
\end{ledgertable}

\note{« Dead? » is written above the line at the end of the heading, in the
same hand and a later ink. The asterisk is hers and is not answered on this
leaf.}

\hand{jo}{\emph{Brown Robertson} \quad 415 Madison Ave.}

\begin{ledgertable}{Y{0.16}Y{0.07}Y{0.40}Y{0.08}}{ & & & }
Jan. 31, 22 & 18 & Evening Wind & 12 \\
 & & '' \quad '' & 12 \\
 & & '' \quad '' & 12 \\
 & 15 & House Tops & 10 \\
 & 18 & East Side Interior - gift to McBride. & --- \\
\end{ledgertable}

\hand{jo}{\emph{Robertson \uncertain{Deschamps} Co.} \quad 415 Mad. Ave.}

\begin{ledgertable}{Y{0.16}Y{0.07}Y{0.28}Y{0.28}}{ & & & }
Sept. 12, 23 & 18 & Mohegan Boat & \\
 & 25 & East Side Interior & Returned June 16, 28. \\
 & 25 & Evening Wind & \\
 & 18 & Railroad & \\
\end{ledgertable}

\note{a single brace gathers the four rows and « Returned June 16, 28. » is
written once against it. The dealer's name is offered as read and doubted: the
letters after « De » will support « Deschamps » and will support « Deachamps »,
and the leaf does not settle which. The account is the last on the leaf; the
lower third is blank.}

\section{Leaf 90: Frank K. M. Rehn --- the watercolours, and the etchings}

\sheet{17444}{90}

\hand{jo}{\emph{Frank K. M. Rehn} - 693 - 5\add{th} Ave.}

\hand{jo}{\emph{Watercolors.}}

\begin{ledgertable}{Y{0.13}Y{0.06}Y{0.34}Y{0.09}Y{0.07}Y{0.11}}{ & & & Sold & Amt. & Rec'd}
Oct. 6, 1924 & 150 & Dead Trees. Exhib. 1924. \quad Loaned to New Gallery, Macbeth & & & \\
 & 150 & Parkhurst's House (the Capt.'s House?) & '' & & \\
 & 150 & Houses + Dory (Italian Quarter) - & Oct. 1924 & 100 & J. N. Spalding, Boston \\
 & '' & Haskell's House (terrasses) \add{side view, no} & Oct. 17, 1926 - 133.33 & & Mrs. Coburn of Chicago \\
 & '' & Italian Quarter - 2 women talking - big houses - & 1924 & 100 & Mrs. J. O. Blanchard \\
 & '' & House at the Fort - cold, on rock - & 1924 & 100 & J. N. Spalding \\
 & '' & Beam trawler, the Seal (profile) & & & \\
 & '' & House with a Bay Window - little yellow, cutie house & 1924 & 100 & Arthur Fowler \\
 & '' & Houses of Squam Light - rep. in Arts Nov. 24 & 1924 & 100 & J. N. Spalding, Boston \\
 & '' & the Pine Tree + house. small & 1924 & 100 & \\
 & '' & Deck of a Beam trawler (from stern, red windlass drums) & 1924 & 100 & J. N. Spalding, Mass \\
Oct. \ill{}, 24 & 150 & Houses in Italian Quarter - big yellow + red house - & 1924 & 100 & Geo. Bellows \\
 & & Bow of Beam Trawler - ventilator + cat head. tanks beyond. & & & \\
 & & Portuguese Church - 2 towers over fence, school yard, flag. & & & \\
 & & Shacks at Sanesville & 1924 & 100 & \\
 & & New England House - no sky, house low among rocks & '' & 100 & \\
 & & House o Harbour - big house & '' & 100 & Mrs. J. O. Blanchard \\
 & & Funnel of Trawler - upper deck, from stern, life boats & '' & 100 & \\
 & & Trawler in Dock, \add{bow} seen from dock, hogsheads on dock. & & & \\
 & & Portuguese Quarter - wash - & '' & 100 & Mrs. J. O. Blanchard \\
 & & Apple Trees - small, fence across foreground & & & \\
 & & House with Fence - pale. stoop on house. & & & \\
 & & Eastern Point Light - lighthouse & & & \\
 & & Gloucester Mansion - lamp post \struck{Haskell's - front view} \add{front view - terrasses - in Arts Nov. 1924.} & 1924 & 100 & Geo. Bellows \\
 & & House on the Shore - Italian women in yellow, grey stones & 1924 & 100 & New Orleans Museum \add{Neuman} \\
 & & Rocks at the Fort. dogs, man on hill \quad Sent to Mrs. Coburn, Chicago, on approval, May '26. & & & \\
\end{ledgertable}

\note{« Watercolors. » is underlined in red pencil and heads the block. « Sold
\quad Amt. \quad Rec'd » are written above the ruling at the right; the date,
price and title columns are unheaded. « Dead Trees » and the two rows below it
share one price of 150, and from « Haskell's House » down the price column
carries a ditto mark and then nothing at all. The day of the second October
date is written over itself and cannot be read. The Haskell's House row runs
its sale into the Sold column and out of it - « Oct. 17, 1926 - 133.33 » is
written across the ruling in a heavier ink, with « Mrs. Coburn of Chicago »
boxed beside it. « Sanesville » is as written.}

\hand{jo}{\emph{Etchings}}

\begin{ledgertable}{Y{0.16}Y{0.07}Y{0.28}Y{0.28}}{ & & & }
Nov. 23, 21. & & Night in Park (18) & \\
'' & & Les Deux Pigeons (18) & \\
'' & \$22 & Evening Wind & Returned \\
Nov. 20, \struck{22} & 25 & East Side Interior & \\
'' & 18 & Cow + Rocks & \\
'' & 18 & the Railroad & \\
'' & 15 & House tops & \\
\end{ledgertable}

\hand{jo}{Oct. 8, 1924\\
Out of another collection (returned later,\\
Evening Wind was\\
Sold to Mr. Clark - June 11, 25 \quad \$ 14.67}

\note{« Etchings » is written small at the left, underlined, with an ink blot
beside it. One large brace gathers all seven rows against the single word
« Returned »; a second brace to the right of that gathers them against the four
lines above, which are written across the ruling and are not part of the table.
The opening parenthesis of « (returned later, » is never closed. The year of
the fourth row is written over itself. This block is continued on leaf 95,
which says so.}

\section{Leaf 91: Rehn continued --- the south-west watercolours, the Paris
caricatures, the oils}

\sheet{17125}{91}

\hand{jo}{Oct. 15, 1925 - summer in S. W. \quad \emph{Watercolors}}

\begin{ledgertable}{Y{0.07}Y{0.44}Y{0.11}Y{0.07}Y{0.10}}{ & & & & }
200 - & Poplars (+ adobe hut) blue sky. - Bgt. by Alex. Baldwin, July 1926 (200) & & & 1.33. \\
'' & St. Francis Towers - bgt. by Duncan Phillips for 150. \quad 1925 & Rec'd & & 100 \\
'' & Ranch House - meat drying & & & \\
'' & La Penitente & & & 133.33 \\
'' & D. \& G. Locomotive & bgt. by Mrs. John Osgood Blanchard \quad Dec. 17, 25. & & 133.33 \\
'' & Adobe Houses & & --- & 133.33 \\
'' & Interior (J. in short tail) \quad for Chicago Museum \quad Mrs. Coburn - fall 1925 & Mar. 31, 26 & Rec'd & 133.33 \\
'' & Manhattan Bridge - Fog Museum & & & \\
'' & Self Portrait & & & \\
\end{ledgertable}

\note{« Watercolors » is underlined in red pencil. A brace gathers La
Penitente, the Locomotive and Adobe Houses against « bgt. by Mrs. John Osgood
Blanchard ». On the Locomotive line « Rio » and « grande » are written small
above the line, over the ampersand and over the G; the line as written on the
leaf is « D. \& G. Locomotive ». The figure 1.33 against Poplars is as written
and is not 133.}

\hand{jo}{\emph{Paris \uncertain{Caricatures}}\\
1925 spring?}

\begin{ledgertable}{Y{0.44}Y{0.14}Y{0.10}Y{0.12}}{ & & & }
La Pierreuse - big bust, front, pompadour. Sent to Mrs. Coburn on approval May, 1926 & Oct. 17, 1926 (100) & 66.67 & \\
Type de Belleville - awful with baby & & & \\
Sargent de Ville & & & \\
Le Militaire & & & \\
La Concierge & & & \\
Le Piou - piou - little sentinel with bayonette & & & \\
Fille de Joie & & & \\
La Grisette & & & \\
\uncertain{Le Mereosuén} - Sent to Mrs. Coburn on approval May 1926 & Oct. 17, 1926 (100) & 66.67 & Mrs. Coburn \\
\end{ledgertable}

\note{the heading is underlined in red pencil and the word after « Paris » is
written over itself in a heavier ink; the reading offered cannot be checked
letter by letter. « 1925 spring? » is hers, question mark and all. The last
title is offered as read and doubted: the letters are clear and the name is
not. « bayonette » is as written.}

\hand{jo}{\emph{Oils} - Feb. 25, 25 \quad Bootleggers \quad 750\\
Fall of 1925 \quad Shown at Bklyn. Museum -\\
Feb. 1926. \quad Tri-National\\
April '' \quad - St. Botolph's Club, Boston}

\begin{ledgertable}{Y{0.14}Y{0.07}Y{0.36}Y{0.20}Y{0.08}}{ & & & & }
Fall 1925 & 750 & House by a Railroad - invited by New Society - 600 - 200 Com. & To Mr. Steven Clark - Mar. 20, 26 paid June 11, 26 & 400. \\
'' & 750 & N. Y. Pavements (nurse + perambulator & '' & \\
Feb. 1926 & 850 & Hoboken façade (``Sunday'') ``To-Day in American Art'' - Mar. 23, 26 & Nov. 13, 26 - (600) Duncan Phillips & 400 \\
 & & Reproduced in Bklyn. Eagel Mar. 7. 26; N. Y. American, Mar. 21, 26, + Phillips memorial catalogue. & & \\
\end{ledgertable}

\hand{jo}{June 2, 26 - \quad \emph{Watercolors}}

\begin{ledgertable}{Y{0.44}Y{0.30}}{ & }
Skylights at 3 Wash. Sq. N. & \\
Chimneys from \add{Roof of Wash. Sq.} '' \quad '' & \\
The Lilly Apt. from Man. Bridge & \\
\end{ledgertable}

\note{« Oils » and the second « Watercolors » are underlined in red pencil.
« Steven Clark » and « Eagel » are as written. The brace on the House by a
Railroad row gathers it with N. Y. Pavements beneath, and the parenthesis after
« nurse + perambulator » is never closed. The leaf ends here and the last
quarter is blank.}

\section{Leaves 92 and 93: what went to Rehn from Gloucester and Rockland,
1926}

\sheet{16304}{92}

\begin{ledgertable}{Y{0.18}Y{0.46}Y{0.10}Y{0.11}}{ & & Rec'd & }
Aug. 15'' \add{1926} Sent from Gloucester & Talbot's House (Rockland) - fine white mansard \quad Bgt. by Rehn - Oct. 200 - 1/3 = & 133 1/3 & Nov. 1926. \\
 & Mrs. Achorne's Parlor \quad not framed & & \\
 & Lime Rock Quarry II (Rockland) & & \\
 & Bow of Beam Trawler Widgeon - black smoke stack forward, shore + grey sky & & \\
 & Beam Trawler Seal - bold profile \quad Bgt. by Rehn - Oct. 200 - 1/3. \quad (200) & 133 1/3 & Nov. 1926 \\
 & Bow of Osprey - wooly bear chafing gear - dick - rust color & & \\
Sent from Gloucester Sept. 30'', 1926. & Schooner's hull - flat big floor - buildings at side + surrounding - Rockland - crowded - & & \\
 & Civil War Camp Ground - telegraph pole, straw colored fields. & & \\
 & Lime Rock Quarry I & & \\
 & Trawler (profile, below dock) + telegraph pole, spikes to climb & & \\
 & Rockland Harbor - bright red house, white house, blue water + sky + pile of bricks \add{pregnant} & & \\
 & Shipyard, Rockland - long side of boat + rows. Boat the Butterfly & & \\
 & Schooner's Bowsprit - derelict drawn ashore, gloomy looking. & & \\
 & Lime Rock R. R. - White house (shingles) + dark tree, tracks across foreground & & \\
 & Beam trawler Osprey - delicate profile, delicate \add{light} wharf floor foreground. & & \\
 & Beam trawler Widgeon - chocolate profile + rigging - far off shore sharp - & & \\
 & Haunted House - old boarded up boarding house, Rockland, near shipyard. & & \\
 & Trees. East Gloucester - wild + wooly & & \\
\end{ledgertable}

\note{« Rec'd » is the only heading on the leaf; the other three columns are
unheaded and are left so. The two dates in the left column are written across
two and three lines each and are given here on the row they open. « 1926 » is
written small above « Aug. 15'' ». « dick » and « wooly » are as written. The
raised fractions of « 133 1/3 » and « 200 - 1/3 » are set on the leaf as a
numerator above a denominator and are written out here.}

\sheet{17434}{93}

\begin{ledgertable}{Y{0.20}Y{0.60}}{ & }
(cont.) \quad Nov. \ill{} & Universalist Church - tower + roof gable pregnound. \\
 & the Hill. Gloucester St. looking down - shadows crossing st. \\
Feb. 7, 27 & Tony's House - bright blue sky. (where rosetta lived) got ink spot, stone wall, happy on \struck{Italian quarter} \\
Feb. 7. 1927 & Abbot's House - Oct. cold day \\
 & House by Squam River - from rear - light \\
Nov. 1\ill{} \quad Sent Aug. 15, 26 \quad from Gloucester & Davis House - Middle St. - yellow sun light on house + tree. \\
 & Anderson's House - dark ish. strong. \\
Painted spring 1926 Brought to R. later in fall - not so good - & Manhattan Bridge entrance - big arch - best of the Lot. \\
 & Harbor Shore. Rockland - beach, rocks, house above \\
'' \quad '' & R. R. Crossing - empty foreground, Rockland. \\
sent - Aug. 15'', 26 & Gloucester Harbor - rocks, breakwater effect, sharp \quad far shore in detail, not very unique composition \\
\end{ledgertable}

\note{this leaf continues leaf 92 without repeating its heading, and says so:
« (cont.) » opens the left column. The left column is the same unheaded date
column as on leaf 92, and here it carries judgements as well as dates -
« Painted spring 1926, brought to R. later in fall - not so good ». Those
notes are written across two, three and four lines each and are given on the
row they stand against; « not so good » is underlined on the leaf. The same
notes are visible from the other side in the photograph of leaf 92 and are not
repeated there. A small cross is drawn between several of the rows. Two dates
in the left column, both beginning « Nov. », run off into the gutter and what
follows the month is lost; the « F » of the two February dates is cut the same
way in this photograph but is legible in the photograph of leaf 92, which
catches this column from the other side of the opening. « pregnound » is as
written. The lower two thirds of the leaf is blank.}

\section{Leaf 94: four oils ready for Rehn}

\sheet{18195}{94}

\hand{jo}{Oils - ready for Hopper}

\hand{jo}{\add{Automat}}

\begin{ledgertable}{Y{0.28}Y{0.10}Y{0.44}}{ & & }
The \struck{Cafeteria} & 28 x 36 & Mr. Schaeffer - \quad girl in green coat, yellow hat at white table, radiator glass window, night outside, day in. \\
The City & 28 x 36 & Houses around a park. Corner house - Duncan Phillips \quad mansard type, many windows, roves + side walls. \\
11 A. M. & 28 x 36 & Nude in arm \struck{chair} chair at window, - Sold Ap. 21. \quad light thru shaft, table, big lamp shade. \quad Man who owns du Bois pictures \\
Gloucester St. & 28 x 36 & 3 gables, very blue sky. \\
Two on the Isle - scrap. & & \\
\end{ledgertable}

\note{« Automat » is written above the ruling, over the struck title of the
first entry, and is not part of the table; « The Cafeteria » is struck through
beneath it. « Oils - ready for Hopper » stands sideways in the left column
across the first two rows. A single diagonal ink stroke runs down the whole
block. The leaf carries nothing below these five lines and two thirds of it are
blank.}

\section{Leaf 95: the etchings continued from leaf 90}

\sheet{16382}{95}

\hand{jo}{Etchings continued from p. 90.}

\begin{ledgertable}{Y{0.15}Y{0.06}Y{0.44}Y{0.11}}{ & & & }
Mar. 15, 26 & 22 & Aux Fortifications & \\
for St. Bot. Cl. show & \$18 & Amer. Landscape & \\
2'' copy Ap. 7, 1926 & 18 & Mohegan Boat - Sold at St. Botolph's Cl., Ap. 1926 - 2 copies * & June 15. 1926. \\
 & 18 & Cat Boat & \\
 & 18 & the Lonely House & \\
 & 18 & House by a River & \\
 & 18 & Cow + Rocks & \\
 & 25 & Night Shadows & \\
 & 18 & the R. R. & \\
 & 15 & Girl on a Bridge & \\
 & 25 & Night in the Park & \\
 & 25 & East Side Interior - Sold at St. Botolph's Cl. - April 26 * & June 15 1926 \\
Nov. 24, 26 & 25 & '' \quad '' \quad '' & \\
 & 25 & Evening Wind - Sold at St. Botolph's Club - Ap. 1926 * & June 15 1926 \\
Jan. 17, 26 & \struck{25} & '' \quad '' (sold copy replaced) & \\
 & 15 & House Tops & \\
 & 18 & Locomotive - Sold at St. Botolph's Cl. Ap. 1926 * & June 15 1926 \\
\struck{\ill{}} & & Les Deux Pigeons '' \quad '' at St. Botolph's Cl - \struck{5} 6 Etchings - \struck{\ill{}} \quad Rec'd 81.34 & June 15 1926 \\
Nov. 24, 26 & & Evening Wind & \\
'' \quad '' \quad '' & & Girl on a Bridge & \\
'' \quad '' \quad '' & & Night in the Park. & \\
More Mar. 15 1926 & & Night - L Train & \\
 & & R. R. Crossing & \\
 & & Train + Bathers & \\
 & & House Tops & \\
 & & Maine Coast & \\
 & & Les Deux Pigeons & \\
 & & Bull Fight & \\
Feb. 3, 1927 & & Aux Fortification - framed & \\
 & & East Side Interior \quad '' & \\
 & & The Locomotive \quad '' & \\
 & & The Cat Boat \quad '' & \\
 & & the R. R. \quad '' & \\
Feb. 8, 27 & & Amer. Landscape \quad '' & \\
 & & Evening Wind \quad '' & \\
 & & Night Shadows \quad '' & \\
 & & Night in the Park \quad '' & \\
 & & Cow + Rocks & \\
\end{ledgertable}

\note{no column on this leaf is headed. The asterisks against the five
St. Botolph's sales are hers and nothing on the leaf answers them. On the Les
Deux Pigeons line the number of etchings is a 6 written over a struck 5, and a
struck word between « Etchings - » and « Rec'd » is unreadable; the date in the
left column of that row is struck through and is unreadable too. The dates in
the left column are as written, and « Jan. 17, 26 » stands below « Nov. 24, 26 »
on the leaf. A single vertical ink stroke runs the length of the block. The last
row is struck through with a horizontal line.}

\section{Leaves 96 to 98: J. H. in the rear studio}

\sheet{17061}{96}

\hand{jo}{J. H. in rear studio\\
3 Wash. Sq. N.}

\note{the leaf is written in pencil in two parallel columns which are not row
for row against one another; they are given here one after the other, left
first. Neither column is ruled off from the other and neither is headed except
where noted.}

\hand{jo}{\emph{Etchings}\\
\quad Lion of \# 6 Wash. Sq. N.\\
\quad Washington Sq. winter \quad 36 x 24\\
\emph{Canvases of J. H.}\\
\quad Head of Frances Fernandez\\
\quad Pastel - Little Ballerina - Marg. Van Epps\\
\quad \quad Kuehne frame \quad at 11.\\
\quad Bodies of trees - Big willow \struck{\ill{}} at Roots\\
\quad Hidalgo Plaza - Monterrey\\
\quad \uncertain{Arabe} patio - \uncertain{Sorillo}\\
\quad Head of Eve Roscoe\\
\quad '' \quad '' Gladys Meyer}

\hand{jo}{\emph{Water Colors framed}\\
\quad Boys on sheet iron crate \quad Quarles\\
\quad John + Jane in boat.}

\note{the words offered and doubted are legible on the leaf and unidentifiable
as names; the struck word between « Big willow » and « at Roots » cannot be
read at all. « Kuehne frame » is written as one run. The right-hand column is
much the shorter of the two and stops half way down the leaf.}

\sheet{16956}{97}

\hand{jo}{\emph{Water Colors of Jo N. Hopper.}}

\hand{jo}{\emph{at Whitney Club} - Nov. 1927. in portfolio}

\begin{ledgertable}{Y{0.44}Y{0.09}Y{0.28}}{ & & }
Lamp in window - shown Dec. 1927. & 50 & \\
Woodstock - Kiddy party & '' & \\
Guinney fleet in fog. & '' & \\
Petunias - purple & '' & Gone to Morton Gallery Mar. 1928 \\
Zineas - tall old one & '' & \\
Moxie Theatre. & 35 & \\
Falls Village - shack, purple + magenta. & 50 & \\
Children come to see the Lady paint & '' & \\
Silver mat + lot of loose flowers, scrapers & 10 x 5. & Taken home. \\
\end{ledgertable}

\hand{jo}{\emph{Members Annual Show - Whitney Club -}}

\begin{ledgertable}{Y{0.55}Y{0.09}}{ & }
Charlestown Tree - E. H.'s mat + frame. & 75. \\
\end{ledgertable}

\hand{jo}{\emph{Morton Gallery} - \quad Mar. 1928.}

\begin{ledgertable}{Y{0.10}Y{0.55}Y{0.09}}{ & & }
White frames \} & Trees in Chase's yard - Woodstock 1921 - tree flaming up, yellow & 50 \\
 & Hill from Athens - \add{on} cheap paper - shown at Bklyn. Mus. + Art Inst. 1923 & 50 \\
 & Petunias - gilt edge mat - very purple - in silhouette \quad for flower show in April. & 50 \\
\end{ledgertable}

\hand{jo}{See page 62.}

\hand{jo}{Way back \quad Chas. Daniels - small oils in white frames. \quad
Exhibited in show there}

\hand{jo}{\emph{Later - the New Gallery} - mostly modern french.}

\hand{jo}{\emph{Water Colors in Exhibitions.}}

\begin{ledgertable}{Y{0.13}Y{0.55}Y{0.13}}{ & & }
1922 ± & Children in Apple tree - sold to Mrs. Francis Bayless Clark. & \\
 & Penn. Academy \quad Glacial Foundations & \\
 & Golden Gate San Francisco Fair - Chez Hopper I. & 1939 or 40 \\
 & Chicago \uncertain{Art} \struck{Watercolor} \struck{Corcoran} Water Color show - Holy City, Wyoming. & \\
 & Bertha Sheafer Gal. - May 1945 - Little stove. & \\
 & Corcoran - S. Truro Hills (+ back your house) oil. & 1943. \\
 & \uncertain{Wythe} - Cat Show - Sleeping cat water col\ill{} & 1942 ± Dec. \\
\end{ledgertable}

\hand{jo}{Mar. 18'' '46. Turned down by jury of \ill{}ush Club (thy lowest
depths \ill{} Pen \& Brush!\\
1947 - Early spring Cat Show - Feragil.. Cats in \ill{} the portfolio.\\
\struck{Exhibited} '' \quad '' \quad '' \quad \uncertain{Scarpino} co\ill{} at
the Wake cat show -}

\note{the fore-edge of the leaf is torn away at the foot and a hole is torn
through the last three lines, taking words out of each; what is marked
illegible there is missing paper, not a hard reading. The fragment « ...ush
Club » is all that survives of the name of the club whose jury refused her, and
it is not completed here from the « Pen \& Brush! » at the end of the same
line, which is a separate exclamation. « Feragil » and « Zineas » are as
written. « Chez Hopper I » is as written, numeral and all. The heading « Later
- the New Gallery » is underlined; the two lines about Chas. Daniels are not
part of any table and are written across the leaf.}

\sheet{17161}{98}

\hand{jo}{J. H. Continued\\
San Esteban (24 x 36 oil) \quad Corcoran Biennial. \add{Early} Spring 1951.\\
Interior of the Ascensions - water color - shown in parish house \add{of} the
Ascension - \quad late Lent 1952.}

\note{three lines, written above and across the top rule of the leaf, and the
rest of it is blank. They are the last writing in the volume. « the
Ascensions » and « the Ascension » are as written, in that order.}

\section{The back board}

\sheet{17296}{}

\note{the back board carries the stationer's printed design and no writing in
any hand. The blue 98 legible at its head is printed on the leaf beneath,
showing through the torn upper corner, and the marks at the foot are the
reverse of leaf 98 showing through a second tear; nothing of either can be
read. It is recorded here because it is the last sheet of the volume and
because the front board, which does carry writing, is transcribed at the head
of batch 1.}

\keywords{medium:etchings, medium:watercolors, medium:oils, medium:pastel,
dealer:Frank K. M. Rehn, dealer:E. \& A. Milch, dealer:Babcock Art Galleries,
dealer:Brown Robertson, dealer:Sidney Phillips, dealer:Print Rooms San
Francisco, dealer:Carl Smalley, dealer:Morton Gallery, dealer:Bertha Schaefer,
society:American Society of Graphic Arts, society:Whitney Studio Club,
society:St. Botolph's Club, society:Pen \& Brush Club,
collection:Corcoran Gallery of Art, collection:Pennsylvania Academy,
person:Josephine Nivison Hopper, person:Mrs. Albert Sterner,
person:Mrs. Lewis Larned Coburn, person:Duncan Phillips, person:George Bellows,
person:J. N. Spalding, person:Stephen Clark, place:Gloucester, place:Rockland,
place:Washington Square, place:Monterrey, feature:consignment accounts,
feature:commission, feature:prints returned, feature:price changes,
feature:picture descriptions, feature:work by Jo Hopper, feature:torn leaf}

\end{document}
